\documentclass[12pt]{article}

%--- Packages ----------------------------------------------------------------
\usepackage[margin=1in]{geometry}
\usepackage{amsmath, amssymb}
\usepackage{booktabs}
\usepackage{graphicx}
\usepackage{float}
\usepackage{hyperref}
\usepackage{natbib}
\usepackage{setspace}
\usepackage{caption}
\usepackage{subcaption}

\onehalfspacing

%--- Title -------------------------------------------------------------------
\title{\textbf{Constructing an Immigration Shift-Share Instrument} \\
       \large Project Exercises --- Path 2}
\author{Jessica Ding}
\date{May 4, 2026}

\begin{document}
\maketitle

%=============================================================================
\section{Introduction}
%=============================================================================

This project constructs a shift-share (Bartik) instrumental variable for
immigration across U.S.\ commuting zones using five decades of Census
microdata (1970--2010), following the approach of
\citet{kim2024measuring}. The instrument is then used to estimate the
causal effect of immigration on local wages via two-stage least squares
(2SLS). The full pipeline---from raw IPUMS microdata to regression
output---is implemented in a single Stata do-file and is fully
reproducible given the raw data.

%=============================================================================
\section{Data}
%=============================================================================

\subsection{Sources}

Individual-level microdata come from the U.S.\ Decennial Census and
American Community Survey, accessed through IPUMS USA
\citep{ruggles2024ipums} for five decades: 1970, 1980, 1990, 2000, and
2010. All five samples were downloaded via the IPUMS API using
\texttt{ipumspy}. Geography is harmonized to 1990-definition commuting
zones using crosswalks from \citet{dorn2009essays}. Birthplace codes are
mapped to approximately 82 origin-country groups using the crosswalk
provided by \citet{kim2024measuring}.

\subsection{Sample Restrictions}

For each decade I impose the following restrictions: (1) drop Alaska and
Hawaii (\texttt{statefip} $\in \{2, 15\}$); (2) drop individuals in
non-standard group quarters (\texttt{GQ} $\in \{0, 3, 4\}$); (3) drop
individuals under age 16. Geography is converted to commuting zones using
a many-to-many join (\texttt{joinby}), with person weights adjusted by
the allocation factor supplied in each Dorn crosswalk
($\tilde{w}_i = \text{afact} \times \text{perwt}$).

\subsection{Wage Construction}

Hourly wages are computed as annual wage income divided by annual hours
worked. Annual hours are the product of weeks worked and usual hours per
week. For 1970 and 1980, weeks worked (\texttt{wkswork2}) is recorded in
intervals; I convert to continuous values using midpoint imputation
(e.g., the interval ``40--47 weeks'' is assigned 43.5 weeks). Wage
observations at the top 1\% of the weighted distribution are dropped to
limit the influence of top-coded incomes.

%=============================================================================
\section{Empirical Strategy}
%=============================================================================

\subsection{Shift-Share Instrument}

Local immigration is endogenous to local labor demand: immigrants may
sort toward booming areas, biasing OLS estimates of the wage effect
upward (or downward in a labor-demand framework). Following
\citet{card2001immigrant} and \citet{kim2024measuring}, I construct a
shift-share (Bartik) instrument that predicts local immigration from
\emph{national} inflows by origin country, distributed across commuting
zones using \emph{historical} (1970) settlement patterns.

The instrument for commuting zone $c$ in decade $t$ is:
\begin{equation}
    \widetilde{B}_{c,t} = \frac{1}{\text{pop}_{c,1970}}
    \sum_{o} \underbrace{\frac{\text{imm}_{c,o,1970}}
    {\text{imm}^{\text{nat}}_{o,1970}}}_{\text{1970 share}}
    \times
    \underbrace{\left(\text{imm}^{\text{nat}}_{o,t} -
    \text{imm}^{\text{nat}}_{o,t-1}\right)}_{\text{national shift}}
    \label{eq:bartik}
\end{equation}
where $\text{imm}_{c,o,1970}$ is the number of immigrants from origin
$o$ living in commuting zone $c$ in 1970, $\text{imm}^{\text{nat}}_{o,1970}$
is the total U.S.\ immigrant population from origin $o$ in 1970, and
$\text{pop}_{c,1970}$ is the total 1970 population of commuting zone
$c$.

\subsection{Identification}

The instrument is valid if the 1970 settlement shares are uncorrelated
with decade-$t$ local labor-demand shocks, conditional on controls. This
is plausible because immigrants' location choices in 1970 reflect
historical network formation rather than current economic conditions.
Identification follows the framework of
\citet{goldsmith2020bartik}.

\subsection{Estimating Equations}

\paragraph{First stage.}
\begin{equation}
    \Delta \text{imm}_{c,t} = \beta\,\widetilde{B}_{c,t}
    + \gamma_t + \varepsilon_{c,t}
\end{equation}
where $\Delta \text{imm}_{c,t} = (\text{imm}_{c,t} -
\text{imm}_{c,t-1})/\text{pop}_{c,1970}$ is the actual immigration
change normalized by 1970 population, and $\gamma_t$ are decade fixed
effects. Standard errors are clustered at the commuting-zone level.

\paragraph{Wage regression.}
\begin{equation}
    \Delta \text{wage}_{c,t} = \alpha\,\Delta \text{imm}_{c,t}
    + \gamma_t + u_{c,t}
\end{equation}
I estimate this equation by OLS and by IV/2SLS, instrumenting
$\Delta\text{imm}_{c,t}$ with $\widetilde{B}_{c,t}$. Both specifications
include year fixed effects and CZ-clustered standard errors. I also
report specifications that add commuting-zone fixed effects.

%=============================================================================
\section{Results}
%=============================================================================

\subsection{Summary Statistics}

Table~\ref{tab:summary} reports means and standard deviations of the
key variables by decade. The shift-share instrument and actual
immigration change track each other closely, consistent with a strong
first stage. Both grow over the sample period, reflecting rising
immigration to the United States from 1970 to 2000, with a modest
decline in the 2010 decade.

% --- paste tabstat output here or use esttab/estout to auto-generate ---
% \input{../output/tables/summary.tex}

\subsection{First Stage}

Figure~\ref{fig:fs_scatter} plots the shift-share instrument against the
actual immigration change for each decade. The positive relationship is
visible in all four panels, with a steeper slope in later decades when
immigration was more concentrated in historically immigrant-receiving
commuting zones.

\begin{figure}[H]
    \centering
    \includegraphics[width=\textwidth]{../output/figures/first_stage_scatter.png}
    \caption{First-stage scatter plots by decade. Each point is a
    commuting zone. The fitted line is from OLS.}
    \label{fig:fs_scatter}
\end{figure}

Table~\ref{tab:fs} reports the first-stage regression. The instrument
enters with a large, statistically significant coefficient in all
specifications. The Kleibergen-Paap F-statistic exceeds the conventional
threshold of 10, indicating that the instrument is strong.

\begin{table}[htbp]\centering
\def\sym#1{\ifmmode^{#1}\else\(^{#1}\)\fi}
\caption{First Stage: IV to Actual Immigration Change}
\begin{tabular}{l*{2}{c}}
\hline\hline
            &\multicolumn{1}{c}{(1)}&\multicolumn{1}{c}{(2)}\\
            &\multicolumn{1}{c}{(1) Year FE}&\multicolumn{1}{c}{(2) Year+CZ FE}\\
\hline
imm\_shift   &       0.384\sym{***}&       0.457\sym{***}\\
            &     (0.097)         &     (0.121)         \\
\hline
\(N\)       &        2775         &        2772         \\
CZ FE       &          No         &         Yes         \\
Year FE     &         Yes         &         Yes         \\
\hline\hline
\multicolumn{3}{l}{\footnotesize Standard errors in parentheses}\\
\multicolumn{3}{l}{\footnotesize \sym{*} \(p<0.1\), \sym{**} \(p<0.05\), \sym{***} \(p<0.01\)}\\
\end{tabular}
\end{table}


\subsection{Wage Regressions}

Table~\ref{tab:wages} reports OLS and IV estimates of the effect of
immigration on local wages. The OLS estimate (column 1) is 0.743
(s.e.\ = 0.446), suggesting that a one-unit increase in the immigration
rate is associated with a \$0.74 increase in mean hourly wages. The
IV estimate (column 2) is 0.775 (s.e.\ = 0.896), very close to OLS.
The first-stage F-statistic is 16.15, well above the conventional
threshold of 10, confirming instrument strength.

When commuting-zone fixed effects are added (columns 3 and 4), the
estimates rise substantially. The IV+FE estimate is 4.073 (s.e.\ = 2.392),
significant at the 10\% level, with a first-stage F of 8.23.

The similarity between OLS and IV in the baseline specification suggests
limited endogeneity bias once year fixed effects absorb aggregate demand
variation. The larger IV estimate with CZ fixed effects is consistent
with immigrants sorting into slower-growing commuting zones within the
panel: controlling for CZ-level trends reveals a larger positive wage
effect of immigration.

\begin{table}[htbp]\centering
\def\sym#1{\ifmmode^{#1}\else\(^{#1}\)\fi}
\caption{Effect of Immigration on Wages}
\begin{tabular}{l*{4}{c}}
\hline\hline
            &\multicolumn{1}{c}{(1)}&\multicolumn{1}{c}{(2)}&\multicolumn{1}{c}{(3)}&\multicolumn{1}{c}{(4)}\\
            &\multicolumn{1}{c}{(1) OLS}&\multicolumn{1}{c}{(2) IV}&\multicolumn{1}{c}{(3) OLS+FE}&\multicolumn{1}{c}{(4) IV+FE}\\
\hline
imm\_actch   &       0.743\sym{*}  &       0.775         &       1.040\sym{**} &                     \\
            &     (0.446)         &     (0.896)         &     (0.472)         &                     \\
[1em]
imm\_r       &                     &                     &                     &       4.073\sym{*}  \\
            &                     &                     &                     &     (2.392)         \\
\hline
\(N\)       &        2075         &        2075         &        2040         &        2040         \\
CZ FE       &          No         &          No         &         Yes         &         Yes         \\
Year FE     &         Yes         &         Yes         &         Yes         &         Yes         \\
First-stage F&           .         &       16.15         &           .         &        8.23         \\
\hline\hline
\multicolumn{5}{l}{\footnotesize Clustered SEs at CZ level. *** p<0.01 ** p<0.05 * p<0.1}\\
\multicolumn{5}{l}{\footnotesize \sym{*} \(p<0.1\), \sym{**} \(p<0.05\), \sym{***} \(p<0.01\)}\\
\end{tabular}
\end{table}


%=============================================================================
\section{What I Learned About the Data}
%=============================================================================

Several features of the data required careful handling. First, the
geography crosswalks differ substantially across decades: 1970 uses
county group identifiers (\texttt{CNTYGP97}), 1980 uses a different
county group scheme (\texttt{CNTYGP98}), and 1990--2010 use PUMAs. Each
crosswalk carries an allocation factor to handle geographic units that
straddle multiple commuting zones. Second, the wage variable required
substantial cleaning: top-coded incomes and zero or missing hours and
weeks account for a large share of observations, particularly in 1970.
Third, the \texttt{joinby} operation---a many-to-many merge---increases
observation counts relative to the raw extract, because some geographic
units overlap multiple commuting zones; using \texttt{rawsum} rather
than \texttt{sum} in the subsequent collapse is essential to recover
correct population totals.

%=============================================================================
\section{AI Use Disclosure}
%=============================================================================

I used Claude (Anthropic) for assistance with Stata code structure,
debugging geography merge issues, and LaTeX formatting. All empirical
decisions, variable construction choices, and interpretations are my
own. A log of AI interactions is available upon request.

%=============================================================================
\bibliographystyle{aer}
\begin{thebibliography}{9}

\bibitem{card2001immigrant}
Card, David. 2001. ``Immigrant Inflows, Native Outflows, and the Local
Labor Market Impacts of Higher Immigration.'' \textit{Journal of Labor
Economics} 19 (1): 22--64.

\bibitem{dorn2009essays}
Dorn, David. 2009. ``Essays on Inequality, Spatial Interaction, and the
Demand for Skills.'' PhD dissertation, University of St.\ Gallen.

\bibitem{goldsmith2020bartik}
Goldsmith-Pinkham, Paul, Isaac Sorkin, and Henry Swift. 2020.
``Bartik Instruments: What, When, Why, and How.''
\textit{American Economic Review} 110 (8): 2586--2624.

\bibitem{kim2024measuring}
Kim, Gueyon, Cassandra Merritt, and Giovanni Peri. 2024.
``Measuring and Predicting `New Work' in the United States: The Role
of Local Factors and Global Shocks.'' NBER Working Paper No.\ 32526.

\bibitem{ruggles2024ipums}
Ruggles, Steven, et al. 2024. \textit{IPUMS USA: Version 15.0}.
Minneapolis: IPUMS. \url{https://doi.org/10.18128/D010.V15.0}.

\end{thebibliography}

\end{document}
